\documentclass[11pt]{article}
\usepackage[margin=1in]{geometry}
\usepackage{amsmath,amssymb,amsthm}
\usepackage{graphicx}
\usepackage{microtype}
\usepackage{xcolor}
\usepackage{enumitem}
\usepackage[colorlinks=true,linkcolor=blue!55!black,citecolor=blue!55!black,urlcolor=blue!55!black]{hyperref}

\newtheorem{theorem}{Theorem}
\newtheorem{lemma}[theorem]{Lemma}
\newcommand{\T}{\mathbb T}
\newcommand{\Z}{\mathbb Z}
\newcommand{\E}{\mathbb E}
\newcommand{\orth}{\mathrm{orth}}

\title{The sharp norm asymptotic of the orthogonal cycle projection\\
\large Candidate full answer to the orthogonal branch of Question 32 in arXiv:2305.12582}
\author{Automated Mathematics Run \texttt{fa\_banach\_001}, lane 12}
\date{9 August 2026}

\begin{document}
\maketitle

\begin{center}
\fcolorbox{blue!55!black}{blue!4}{\parbox{0.89\textwidth}{
\textbf{Status.} Candidate full solution of the $P_{\orth}$ branch, subject
to human verification.  If $P_{\orth}$ is the orthogonal projection from
the edge space of the discrete two-dimensional torus onto its cycle space,
then
\[
 \|P_{\orth}\|_{\ell_1(E_n)\to\ell_1(E_n)}
 =\frac4\pi\log n+O(1).
\]
The leading constant is sharp.  The separate $P_{\min}$ branch remains open.}}
\end{center}

\section{Source and exact question}

The source is S.~J.~Dilworth, D.~Kutzarova, and M.~I.~Ostrovskii,
\emph{Cycle spaces: invariant projections and applications to transportation
cost}~\cite{DKO}, arXiv:2305.12582.  For the graph
$\T_n=\Z_n\times\Z_n$, let $E_n$ consist of the positively oriented
horizontal and vertical edges, let $Z_n$ be its cycle space, and let
$P_{\orth}:\ell_2(E_n)\to Z_n$ be the orthogonal projection, regarded also
as an operator on $\ell_1(E_n)$.

After proving only
\[
 a_1\sqrt{\log n}\le \|P_{\min}\|_1\le a_2\log n,
\]
the authors state Question 32 on source PDF page 18:
``Find good asymptotic estimates (improving on Proposition 31) for
$\|P_{\min}\|_1$ and/or $\|P_{\orth}\|_1$.''  Figure~\ref{fig:question}
shows the statement and its immediate context.

\begin{figure}[ht]
 \centering
 \includegraphics[width=\textwidth]{figures/open_problem_crop.png}
 \caption{Question 32 on page 18 of the arXiv source~\cite{DKO}.}
 \label{fig:question}
\end{figure}

The theorem below gives a sharp answer to the orthogonal-projection branch.
It does not identify $P_{\min}$, which the source shows already differs from
$P_{\orth}$ when $n=6$.

\section{Idea of the proof}

The orthogonal complement of the cycle space is the cut, or gradient,
space.  Thus the complementary projection is the discrete Helmholtz
projection
\[
 Q_n=I-P_{\orth}=d\Delta^\dagger d^*,
\]
where $d$ is the forward difference, $\Delta=d^*d$, and the dagger denotes
the inverse on mean-zero vertex functions.  Translation and coordinate
symmetry reduce its $\ell_1$ operator norm to the absolute sum of one
column.  That column is a discrete dipole field: the two relevant kernels
are second differences of the Green function.

At distance $r$ from the pole, the Fourier symbol has the same singular
part as $\xi_j\xi_1/|\xi|^2$.  Its inverse Fourier transform is
\[
 H_{11}(x)=-\frac{\cos(2\theta_x)}{2\pi r^2},\qquad
 H_{21}(x)=-\frac{\sin(2\theta_x)}{2\pi r^2}.
\]
A singular-symbol estimate makes the error $O(r^{-3})$.  Passing from the
infinite lattice to the finite torus periodizes $H$; square symmetry cancels
the degree $-2$ tail, leaving a bounded scaled periodic kernel.  Hence all
errors have bounded total $\ell_1$ mass.

The main term contributes
\[
 \frac1{2\pi}\int_1^n\frac{dr}{r}
 \int_0^{2\pi}(|\cos2\theta|+|\sin2\theta|)\,d\theta
 =\frac4\pi\log n+O(1).
\]
Finally, a trace computation shows that replacing $Q_n$ by $P_{\orth}$
adds exactly $n^{-2}$ to the norm.

\section{Fourier kernel estimates}

Write $e_1=(1,0)$, $e_2=(0,1)$ and
\[
 (df)_j(x)=f(x+e_j)-f(x),\qquad \Delta=d^*d.
\]
For $\vartheta\in(-\pi,\pi]^2$, put
\[
 b_j(\vartheta)=e^{i\vartheta_j}-1,\qquad
 \lambda(\vartheta)=|b_1(\vartheta)|^2+|b_2(\vartheta)|^2,
\]
and, away from the origin,
\[
 m_j(\vartheta)=\frac{b_j(\vartheta)\overline{b_1(\vartheta)}}
 {\lambda(\vartheta)}\qquad(j=1,2).
\]
The value assigned at $0$ will never matter beyond an $O(n^{-2})$ term.

For $x\ne0$, define the homogeneous kernels
\begin{equation}\label{eq:H}
 H_j(x)=\frac1{2\pi}
 \left(\frac{\delta_{j1}}{|x|^2}-\frac{2x_jx_1}{|x|^4}\right).
\end{equation}

\begin{lemma}[Infinite-lattice singular symbol]\label{lem:infinite}
Let
\[
 \widehat m_j(a)=\frac1{(2\pi)^2}\int_{(-\pi,\pi]^2}
 e^{ia\cdot\vartheta}m_j(\vartheta)\,d\vartheta,
 \qquad a\in\Z^2.
\]
Then, uniformly for $a\ne0$,
\[
 \widehat m_j(a)=H_j(a)+O((1+|a|)^{-3}).
\]
\end{lemma}

\begin{proof}
Taylor expansion at the unique singular point gives
\begin{equation}\label{eq:symbol-expansion}
 m_j(\vartheta)=\frac{\vartheta_j\vartheta_1}{|\vartheta|^2}
 +r_j(\vartheta),\qquad
 |\partial^\alpha r_j(\vartheta)|\le C_\alpha
 |\vartheta|^{1-|\alpha|}.
\end{equation}
Choose a smooth cutoff $\chi$ supported near zero and equal to one in a
smaller neighborhood.  The part of $m_j$ away from zero is smooth on the
Fourier torus, so repeated integration by parts makes its coefficients
$O_N(|a|^{-N})$.

The inverse Fourier transform on $\mathbb R^2$ of
$\vartheta_j\vartheta_1/|\vartheta|^2$ is precisely $H_j$: this follows by
differentiating the fundamental solution
$-(2\pi)^{-1}\log|x|$ of the positive Laplacian.  Multiplication by $\chi$
convolves $H_j$ with a Schwartz function of integral one.  Since
$|\nabla H_j(x)|\le C|x|^{-3}$, splitting that convolution into
$|y|\le|a|/2$ and its Schwartz tail shows that the result is
$H_j(a)+O(|a|^{-3})$.

It remains to estimate $\chi r_j$.  A smooth dyadic decomposition into
frequency annuli of radius $\rho$ turns \eqref{eq:symbol-expansion} into
the standard bound
\[
 |\mathcal F^{-1}(\chi r_j)_\rho(a)|
 \le C_N\rho^3(1+\rho|a|)^{-N}.
\]
Summing separately over $\rho\le |a|^{-1}$ and
$\rho>|a|^{-1}$ gives $O(|a|^{-3})$.  Combining the three pieces proves
the assertion.
\end{proof}

For $x\in\T_n$, choose its representative in
$[-n/2,n/2)^2$ and write $\rho_n(x)=|x|$.

\begin{lemma}[Torus dipole kernel]\label{lem:torus}
Define
\begin{equation}\label{eq:Kn}
 K_{j,n}(x)=\frac1{n^2}\sum_{k\in\Z_n^2\setminus\{0\}}
 e^{2\pi i k\cdot x/n}m_j(2\pi k/n).
\end{equation}
There is an absolute constant $C$ such that
\begin{align}
 K_{j,n}(x)&=H_j(x)+O(\rho_n(x)^{-3}+n^{-2})
 &&(1\le\rho_n(x)\le n/8),\label{eq:local}\\
 |K_{j,n}(x)|&\le Cn^{-2}
 &&(\rho_n(x)>n/8).\label{eq:far}
\end{align}
\end{lemma}

\begin{proof}
We include the periodization step because it is the point at which a
non-summable absolute $|a|^{-2}$ estimate would lose the theorem.
Fourier aliasing, first for Abel-regularized Fourier series and then by
letting the regularization parameter tend to zero, gives
\begin{equation}\label{eq:alias}
 K_{j,n}(x)=\sum_{\ell\in\Z^2}^{\rm sym}
 \widehat m_j(x+n\ell)+O(n^{-2}).
\end{equation}
The last term accounts for the arbitrary value of the multiplier at the
single grid point $0$ and for the bounded change from Abel to centered-square
summation of the homogeneous degree $-2$ term.  Lemma~\ref{lem:infinite}
makes the periodization of its remainder absolutely convergent.

For the main term, define on $\mathbb R^2\setminus\Z^2$
\[
 \mathcal H_j(z)=\lim_{N\to\infty}
 \sum_{\|\ell\|_\infty\le N}H_j(z+\ell).
\]
This symmetric limit exists.  Indeed,
$H_1(\ell)=(\ell_2^2-\ell_1^2)/(2\pi|\ell|^4)$, so its sum over every
centered square of nonzero lattice points vanishes under coordinate
interchange.  Likewise $H_2(\ell)=-\ell_1\ell_2/(\pi|\ell|^4)$, whose sum
vanishes under either coordinate reflection.
Consequently,
\[
 \mathcal H_j(z)-H_j(z)
 =\sum_{\ell\ne0}[H_j(z+\ell)-H_j(\ell)],
\]
and this series converges locally uniformly because its tail is
$O(|z||\ell|^{-3})$.  Thus $\mathcal H_j(z)-H_j(z)$ is bounded for
$|z|\le1/8$, while $\mathcal H_j$ itself is bounded on the compact part of
the unit torus at distance at least $1/8$ from the pole.

Homogeneity now gives
\[
 \sum_{\ell}^{\rm sym}H_j(x+n\ell)
 =n^{-2}\mathcal H_j(x/n).
\]
For $|x|\le n/8$, this is $H_j(x)+O(n^{-2})$; away from that disk it is
$O(n^{-2})$.  The absolutely convergent error from
Lemma~\ref{lem:infinite} is $O(|x|^{-3}+n^{-3})$ in the first region and
$O(n^{-3})$ in the second.  Substitution in \eqref{eq:alias} proves
\eqref{eq:local}--\eqref{eq:far}.
\end{proof}

\section{Sharp norm asymptotic}

\begin{theorem}\label{thm:main}
For the two-dimensional discrete torus $\T_n=\Z_n^2$,
\[
 \boxed{\displaystyle
 \|P_{\orth}\|_{\ell_1(E_n)\to\ell_1(E_n)}
 =\frac4\pi\log n+O(1).}
\]
In particular, $\|P_{\orth}\|_1=\Theta(\log n)$.
\end{theorem}

\begin{proof}
The cycle space is $\ker d^*$ and the cut space is $\operatorname{ran}d$.
It follows from the orthogonal decomposition of the edge space that
\[
 Q_n:=I-P_{\orth}=d\Delta^\dagger d^*.
\]
Under the discrete Fourier transform, the $2\times2$ multiplier of $Q_n$
is
\[
 \left(\frac{b_j(\vartheta)\overline{b_\ell(\vartheta)}}
 {\lambda(\vartheta)}\right)_{j,\ell=1}^2
\]
at nonzero frequencies and zero at frequency zero.  Hence, up to harmless
translations caused by the edge orientation, the column of $Q_n$ at the
horizontal edge based at zero has components $K_{1,n}$ and $K_{2,n}$.
All columns have the same absolute sum by translation and coordinate
symmetry, so
\begin{equation}\label{eq:qnorm}
 \|Q_n\|_1=\sum_{x\in\T_n}
 (|K_{1,n}(x)|+|K_{2,n}(x)|).
\end{equation}

The contribution of $x=0$ is bounded because the multiplier matrix is an
orthogonal projection at every frequency.  Lemma~\ref{lem:torus} shows
that the region $\rho_n(x)>n/8$ also contributes $O(1)$, and in the
remaining region the total difference between the summands in
\eqref{eq:qnorm} and $|H_1(x)|+|H_2(x)|$ is bounded by
\[
 C\sum_{1\le|x|\le n/8}(|x|^{-3}+n^{-2})=O(1).
\]
Writing $x=r(\cos\theta,\sin\theta)$ in \eqref{eq:H} gives
\[
 |H_1(x)|+|H_2(x)|
 =\frac{|\cos2\theta|+|\sin2\theta|}{2\pi r^2}.
\]
The sum of this locally Lipschitz homogeneous function over lattice points
differs from its integral over the corresponding union of unit squares by
$O(1)$: the error on a square at radius $r$ is $O(r^{-3})$, and those
errors are summable over annuli.  Therefore
\begin{align*}
 \|Q_n\|_1
 &=\frac1{2\pi}\int_1^{n/8}\frac{dr}{r}
   \int_0^{2\pi}(|\cos2\theta|+|\sin2\theta|)\,d\theta+O(1)\\
 &=\frac1{2\pi}(4+4)\log n+O(1)
 =\frac4\pi\log n+O(1).
\end{align*}

It remains only to pass from $Q_n$ to $P_{\orth}$.  Edge transitivity and
$\operatorname{rank}Q_n=n^2-1$ imply that every diagonal entry of $Q_n$ is
\[
 q_{ee}=\frac{\operatorname{trace}Q_n}{|E_n|}
 =\frac{n^2-1}{2n^2}.
\]
The off-diagonal entries of $P_{\orth}=I-Q_n$ are the negatives of those
of $Q_n$, while the two diagonal entries in the selected columns are
positive.  Thus
\[
 \|P_{\orth}\|_1-\|Q_n\|_1
 =(1-q_{ee})-q_{ee}=\frac1{n^2}.
\]
The asserted formula follows.
\end{proof}

\section{Verification, limitations, and novelty}

\begin{enumerate}[leftmargin=*,itemsep=0.55em]
\item \textbf{Exact finite Fourier check.}
The script \texttt{code/verify\_fft.py} evaluates the multiplier in
\eqref{eq:Kn}.  It reproduces the source values
\[
 \frac{19}{9},\quad\frac{41}{16},\quad\frac{69}{25},\quad
 \frac{3839}{1260}
\]
for $n=3,4,5,6$.  For $n=64,128,256,512$, the residuals
$\|P_{\orth}\|_1-(4/\pi)\log n$ are approximately
$0.71556,0.71351,0.71268,0.71234$.  This computation is not used as proof.

\item \textbf{Sensitive analytic point.}
The leading lattice kernel is not absolutely summable in two dimensions.
Lemma~\ref{lem:torus} therefore uses symmetric/Abel periodization, not a
naive absolute periodization.  The cancellation of $H_j(\ell)$ on centered
squares and the absolutely convergent difference
$H_j(z+\ell)-H_j(\ell)$ are the key checks.

\item \textbf{Scope.}
The result completely determines the first-order asymptotic of
$P_{\orth}$ and hence answers the ``and/or $P_{\orth}$'' branch of
Question 32.  It does not determine $P_{\min}$ or its norm.

\item \textbf{Bounded novelty check.}
The run registry, solution, attempt, and proof-gap indexes were searched for
arXiv:2305.12582 and the core projection/torus terms.  Web searches through
9 August 2026 used the exact title and question, ``orthogonal cycle
projection'', ``discrete torus Green function'', and the proposed
``$4/\pi\log n$'' formula.  No later answer was found.  The published author
offprint still contains the question as Question 32.  This supports, but
does not certify, novelty.

\item \textbf{Human-review recommendation.}
Verify Lemma~\ref{lem:torus}, especially the Abel-regularized aliasing
identity and the uniform boundedness of $\mathcal H_j$ away from its pole.
Then check the Fourier orientation only up to translation; the absolute
column sum is invariant under those shifts.
\end{enumerate}

\begin{thebibliography}{9}
\bibitem{DKO}
S.~J.~Dilworth, D.~Kutzarova, and M.~I.~Ostrovskii,
\emph{Cycle spaces: invariant projections and applications to transportation
cost}, arXiv:2305.12582 (2023), Question 32.
\end{thebibliography}

\end{document}
